% 备赛文档：双方共用。
\documentclass[prep,both]{debate}
\motion{“磕CP”热潮折射出人们对爱情的更深信仰 / 更深怀疑}
\meta{赛制}{招新赛 3v3}
\meta{生成日期}{2026-09-30}
\begin{document}
\debatetitle

\begin{summary}
\begin{fields}
\field{持方优劣评估}{正方略劣。题目的“更深”天然邀请“表面甜、底下苦”的读法，嗑CP“爱情代餐”的说法在舆论里也更常见，评委直觉偏反；正方的论据反而更硬（复旦、中科院两份大样本调查都指向“想要爱情、卡在遇见”）。正方的打法是抢判准的“内容”一半：人们嗑什么暴露了信什么，并把“疑”锁在“遇见”层面；反方的打法是抢判准的“形式”一半：为什么只敢隔着屏幕、结局写好地爱。}
\thesis{正方}{人们嗑的是最严格的爱情标准，疑的只是遇不遇得到；疑在上面，信在底下。}
\thesis{反方}{人们只敢把爱放在屏幕那头、结局写好的地方；信的是别人的爱，疑的是自己的爱靠不靠得住，疑在底下。}
\end{fields}
\end{summary}

\section{一、赛制要点}
\begin{flowtable}
\block{A}
\flow{1 正方一辩立论}{180 秒}{正一}{定义、中立判准、两个论点、全场问题}
\flow{2 反方二辩质询正方一辩}{90 秒 双边}{反二 问}{全场第一次质询：拿“不进入那段关系”的承认，逼“既然信为什么不去谈”}
\flow{3 反方一辩立论}{180 秒}{反一}{开头一句接住质询成果，然后立论主体}
\flow{4 正方二辩质询反方一辩}{90 秒 双边}{正二 问}{拿“嗑的是专一”“疑的是遇见”两个承认}
\block{B}
\flow{5 三辩对辩}{各 90 秒}{正三 与 反三}{开第一个战场：内容与形式，谁更说明心态}
\flow{6 反方二辩申论}{90 秒}{反二}{全场第一篇申论：拆正方立论 + 推进反方论点一}
\flow{7 正方三辩质询反方二辩}{90 秒 双边}{正三 问}{把“怕”变成“当真”，把人数拉回题目主体}
\flow{8 正方二辩申论}{90 秒}{正二}{回应反二申论 + 拆反方立论 + 推进正方论点一}
\flow{9 反方三辩质询正方二辩}{90 秒 双边}{反三 问}{把“标准”变成“不用兑现”，把“宁缺毋滥”变成“缺”}
\block{C}
\flow{10 一辩对辩}{各 90 秒}{正一 与 反一}{收拢前两段：谁是谁的来源}
\flow{11 正方三辩申论}{90 秒}{正三}{处理反方最强点 + 推进正方论点二}
\flow{12 反方一辩质询正方三辩}{90 秒 双边}{反一 问}{为反一结辩取材：人们“不敢”，不敢与不信分不开}
\flow{13 反方三辩申论}{90 秒}{反三}{回应正三申论 + 处理正方最强点 + 推进反方论点二}
\flow{14 正方一辩质询反方三辩}{90 秒 双边}{正一 问}{全场最后一次质询：取“相信好结局存在”“热潮里也有信”}
\block{D}
\flow{15 二辩对辩}{各 90 秒}{正二 与 反二}{结辩前最后交锋：各守判准的一半}
\flow{16 反方一辩结辩}{90 秒}{反一}{先结辩：封路按正方全场主打路线分条件}
\flow{17 正方一辩结辩}{90 秒}{正一}{最后一发：回应反一结辩，约 40 字临场位}
\end{flowtable}
质询 90 秒双边计时，频繁刻意打断要扣分，所以问题写成一句话能答完的封闭式，对方不答先复述一次再礼貌换问。没有驳论、小结和自由辩：对方立论靠两篇申论拆，二辩申论拆立论主干，三辩申论处理最强点。每篇申论之后紧跟被质询，申论里只用已核实的材料。一辩稿算全队作品，一辩的个人分来自被质询、一辩对辩、最后一次质询和结辩。

\section{二、辩题性质与争议焦点}
\begin{fields}
\claim[辩题性质]{事实判断中的比较题：问一个社会现象底下的心态是哪一种。“更深”要求承认两种心态都在，比较谁在底下。}
\field{正方倾向把题目解释成}{看人们在热潮里嗑的是什么样的爱：内容就是标准，标准就是信。}
\field{反方倾向解释成}{看人们用什么方式去爱：隔着屏幕、结局写好、随时可退，方式就是疑。}
\field{争议焦点}{\sub{战场一}{内容与形式，谁更说明心态（正方占内容，反方占形式）。}\sub{战场二}{“不敢”算不算“不信”：疑的是遇见，还是疑爱靠不靠得住。}\sub{战场三}{谁是谁的来源：先有想要才不敢，还是先有不敢才去嗑。}}
\end{fields}

\section{三、关键词定义}
\subsection{3.1 磕CP（嗑CP）}
\begin{fields}
\field{调研}{\sub{词典义}{《现代汉语词典》未收 CP；“嗑（kè）”：用上下门牙咬有壳的或者硬的东西；“磕”：碰在硬的东西上。题面写“磕”，通行写法是“嗑”，场上不纠字（水寒说语文，网易，2022）。}\sub{学术义}{许冠文、张慧（2023）以“嗑CP”为研究对象：青年女性粉丝为角色或人物配对并从中获得积极情感体验（《妇女研究论丛》2023 年第 1 期）。}\sub{通行义}{CP 即 coupling，源自日本同人文化，指观众把人物配成一对并沉迷于两者可能存在的浪漫关系（百度百科“嗑cp”词条，查证用）。}}
\begin{procard}{正方有利定义}\definition{嗑CP，是为一段理想爱情投入情感，并在社群里共同确认它的活动。}强调“投入”和“确认”。\sub{有利之处}{把嗑CP 写成一种积极的情感行为，信仰呼之欲出。}\sub{反方如何反驳它}{把“嗑”写成“追求”，偷换了旁观的本质；嗑的人一步都不用走。}\end{procard}
\begin{concard}{反方有利定义}\definition{嗑CP，是以旁观代替亲历、以想象代替关系的情感消费。}强调“代替”。\sub{有利之处}{替代天然指向“不信正餐吃得到”。}\sub{正方如何反驳它}{“代替”预设了他们本来要去恋爱而没去，是把结论写进定义；许冠文本人说，只停留在替代性满足这一步是不够的。}\end{concard}
\begin{neutralcard}{中立定义}以旁观者的身份，给虚构角色或者真实人物配对，想象并期待他们之间的浪漫关系，并从中得到情感满足。\sub{合理性}{贴近网络通行义与学术研究对象；不预设“代替”，也不预设“追求”，双方都能在里面找到空间。}\end{neutralcard}
\end{fields}

\subsection{3.2 热潮与折射}
\begin{fields}
\field{调研}{\sub{热潮}{嗑CP 从同人小圈子进入大众平台，高校 CP 粉丝群、影视剧 CP、真人 CP 都有大量参与者（许冠文 2023 访谈：晚上发布招募，三小时内加了快 40 个访谈对象）。}\sub{折射}{比喻通过一个现象间接显出另一样东西。}}
\begin{procard}{正方有利定义}\definition{热潮，说明爱情叙事仍是最能调动人的故事。}\sub{有利之处}{规模本身就是信的证据。}\sub{反方如何反驳它}{规模只说明人们爱看，不说明人们信；同一时期婚恋也在退潮。}\end{procard}
\begin{concard}{反方有利定义}\definition{热潮，是与婚恋退潮同时出现的现象，两者要放在一起看。}\sub{有利之处}{把单身率、结婚登记拉进比赛。}\sub{正方如何反驳它}{结婚登记跟着条件走，二〇二五年全国通办后一年就多了 65.7 万对；拿婚恋率量不出爱情观。}\end{concard}
\begin{neutralcard}{中立定义}热潮：嗑CP 从小圈子变成大量普通年轻人都在参与、随处可见的现象。折射：通过这个现象的特征，间接显出参与者的心态，不以口号为准。\sub{合理性}{“我又相信爱情了”是说法，不是折射；双方都要从人们嗑什么、怎么嗑、为什么嗑推出心态。}\end{neutralcard}
\end{fields}

\subsection{3.3 对爱情的信仰 / 怀疑}
\begin{fields}
\field{调研}{\sub{词典义}{信仰：对某种主张、主义、宗教或某人极其相信和尊敬（汉典）；《现代汉语词典》另有“拿来作为自己行动的榜样或指南”（有把握，赛前查纸本）。怀疑：疑惑，不很相信（《现代汉语词典》，有把握）。}\sub{学术义}{浪漫信念：Sprecher 和 Metts（1989）把它拆成四条：爱能找到出路、唯一、理想化、一见钟情。}}
\begin{procard}{正方有利定义}\definition{对爱情的信仰，是相信爱情真实存在、值得，这种相信不以自己亲身得到为前提。}\sub{有利之处}{嗑CP 的人不必去谈，也能算信。}\sub{反方如何反驳它}{不以亲身为前提的相信，就是“别人的爱是真的，我的不会是”，这恰恰是疑；词典说信仰要“作为行动的指南”。}\end{procard}
\begin{concard}{反方有利定义}\definition{对爱情的信仰，是相信爱情会在自己身上发生，并愿意为它承担风险；做不到就是怀疑。}\sub{有利之处}{单身、观望、嗑CP 都能算成疑。}\sub{正方如何反驳它}{这是把“信仰”换成“行动”；按这个标准，所有还没遇到对的人的单身者都不信爱情，等于拿婚恋率判爱情观。}\end{concard}
\begin{neutralcard}{中立定义}对爱情的信仰，是相信真正的爱情存在、值得期待和投入；对爱情的怀疑，是不相信真正的爱情存在，或者不相信它靠得住、值得投入。\sub{合理性}{两个词在同一层面对称。正方主要在“存在、值得”上发力，反方主要在“靠得住、值得投入”上发力，双方都有得打。}\end{neutralcard}
\end{fields}

\subsection{3.4 更深}
\begin{fields}
\begin{procard}{正方有利定义}\definition{更深，是表层是对现实条件的犹豫，底层是对爱情本身的信。}\sub{反方如何反驳它}{把所有犹豫都归到“条件”，等于先把疑定义成浅的。}\end{procard}
\begin{concard}{反方有利定义}\definition{更深，是表层是热闹的甜，底层是对爱情能不能落地的疑。}\sub{正方如何反驳它}{把所有甜都叫“表层”，等于先把信定义成浅的。}\end{concard}
\begin{neutralcard}{中立定义}热潮里两种心态都有，比较哪一种更深，也就是更底层：它是另一种心态的来源，拿掉它，热潮就撑不起来。\sub{合理性}{两版有利定义结构对称，中立版只问“谁在底下”，把证明责任平均分给双方。}\end{neutralcard}
\end{fields}

\subsection{3.5 隐含要素}
\begin{fields}
\field{主体}{参与热潮的人，主要是青年；不是全社会，也不只是极端粉丝。}
\field{范围}{当下中国互联网上的嗑CP 现象。}
\field{时间尺度}{二〇二〇年代。}
\field{比较对象}{同一群人心里的信与疑，比谁在底下；不是跟过去比。}
\end{fields}

\begin{warnbox}{3.6 定义杀陷阱：双方都要背}
\begin{fields}
\begin{procard}{正方的定义杀陷阱}把“信仰”定义成“只要喜欢爱情故事就算信”。评委会觉得反方无路可走，按这个定义看言情小说的人全都信爱情，题目失去意义。\end{procard}
\begin{concard}{反方的定义杀陷阱}把“信仰”定义成“必须亲自去谈、去结婚”。评委会觉得这是拿婚恋率偷换爱情观，而且正方在这个定义下不可能赢。\end{concard}
\end{fields}
\end{warnbox}

\section{四、判准与论证义务}
\begin{fields}
\claim[判准是什么]{哪一种心态更能解释嗑CP 热潮的核心特征：人们嗑的是什么样的爱，为什么用旁观的方式去爱；并且它是另一种心态的来源。}
\begin{procard}{正方有利判准}\definition{看人们在热潮里把什么当成爱情的标准，标准越理想，越说明信。}\sub{有利之处}{只看内容，内容全是专一、双向、好结局。}\sub{反方如何反驳它}{只看内容回避了形式；想象里标准再高也不花成本，叶公对龙的标准也很高。}\end{procard}
\begin{concard}{反方有利判准}\definition{看人们愿不愿意把这份相信用在自己身上，不用在自己身上的信就是疑。}\sub{有利之处}{嗑CP 的形式天然是不用在自己身上。}\sub{正方如何反驳它}{这是把“信仰”换成“行动”；按此判准只要单身率高反方就赢，拿婚恋率判了爱情观。}\end{concard}
\begin{neutralcard}{中立判准}哪一种心态更能解释热潮的核心特征（嗑的是什么、为什么这样嗑），并且是另一种心态的来源。\sub{合理性}{题目说“折射”，就得从现象特征推心态；“更深”要求比较并说明谁是谁的来源；双方义务对称。}\sub{正方需要证明}{嗑CP 的动力来自相信真爱存在、值得；热潮里的犹豫疑的是遇见，不是爱情本身。}\sub{反方需要证明}{嗑CP 的动力来自不信真爱靠得住；热潮里的甜只在安全距离内成立。}\end{neutralcard}
\field{全场问题}{\sub{正方问}{如果人们不信爱情，为什么偏偏嗑“唯一”和“双向奔赴”，而不嗑一对互相算计的人？}\sub{反方问}{如果人们真信爱情，为什么要把它放在一段自己碰不到、结局早写好的关系里？}}
\end{fields}

\prosection{五、正方论点}
\prosubsection{论点一：嗑的是理想}
\begin{fields}
\claim{人们嗑CP，嗑的是最严格的爱情标准：唯一、坚守、双向、好结局。肯为一个还没在现实里见过的标准投入热情，并拿它拒绝将就，这就是信。}
\begin{mechanism}\step 嗑谁、嗑什么是每个人自己挑的，挑什么就暴露心里的标准\step 被挑中的正是浪漫信念的几条：唯一、爱能找到出路、把对方理想化\step 这套标准被带回现实，年轻人拿它拒绝将就\end{mechanism}
\begin{evidence}{学理}[已核实]浪漫信念量表（Sprecher 和 Metts，1989）\sub{核心主张}{浪漫主义可以拆成四条信念：爱能找到出路、唯一、理想化、一见钟情。}\sub{支撑}{机制第二步：CP 叙事几乎就是这四条的演出。}\sub{边界}{量表测的是人对自己恋爱的信念，不直接测对 CP 的期待；只拿来说明“浪漫信念长什么样”，不说它测了 CP。}\sub{口头版}{心理学里的浪漫信念就四条，CP 剧全占了。}\sub{出处}{Journal of Social and Personal Relationships 6(4): 387–411，730 名本科生。}\end{evidence}
\begin{evidence}{数据}[已核实]许冠文、张慧（2023）深访 53 名女性 CP 粉；许冠文说，CP 更像是一个将她们的期待具象化的载体，只停留在“安抚毯”式的替代性满足是不够的。\sub{来源}{许冠文、张慧，《“嗑CP”：青年女粉丝的创造性情感体验》，《妇女研究论丛》2023 年第 1 期；红星新闻 2023-10-14 采访，凤凰网转载。}\sub{方法}{质性深访，经高校 CP 粉丝群招募，55 人中采用 53 名女性，本科或研究生学历。}\sub{局限}{高学历女性、滚雪球招募，不能推到全体；研究者本人也说现实亲密关系“每一步都充满风险”，反方会引后半句。}\end{evidence}
\begin{evidence}{案例}[已核实]“双向奔赴”入选《咬文嚼字》2023 年十大流行语，释义：“多用于人与人之间，表达了人们相互爱慕、相互亲近的美好愿望。”\sub{代表性}{年度流行语是大众语言的样本，不是小圈子用语。}\sub{贴合度}{说明“双向”这一爱情理想被全社会接受；不要说这个词来自 CP 圈，没有出处。}\sub{出处}{文汇报客户端 2023-12-04，中国作家网转载。}\end{evidence}
\begin{evidence}{补充数据}[已核实]中国青年报社会调查中心：对一段恋爱的期待，58.5\% 的受访青年希望互相欣赏、变得更加自信，58.1\% 期待彼此包容、更轻松有安全感。\sub{来源}{杜园春，《年轻人婚恋更看重情感质量》，中国青年报 2024-05-30 第 3 版。}\sub{局限}{文中没有公开样本量与调查方法；只用于正二申论与被质询防守，场上被问样本就如实说“文中未公开”。}\end{evidence}
\closing{内容特征只能用信来解释：真不信爱情的人会嗑算计、嗑利用，不会嗑“双向奔赴”；而且这把尺子被带回了现实。}
\begin{clash}
\pair{标准高是因为不用兑现，叶公对龙的标准也很高。}{叶公怕真龙；年轻人不怕爱，怕将就。标准被带回了现实，“宁缺毋滥”就是在用尺子。}
\pair{量表测的是对自己的信念，你们拿来量 CP，自变量换了。}{我方只用量表说明浪漫信念有哪几条，CP 剧演的正是这几条，这是对照，不是测量。}
\pair{“双向奔赴”后来用在了外交上，不是爱情。}{释义第一句就是“多用于人与人之间，表达相互爱慕的愿望”；用得越广，越说明这个理想被接受。}
\end{clash}
\end{fields}

\prosubsection{论点二：疑的是条件}
\begin{fields}
\claim{热潮里的犹豫是真的，但它疑的是遇不遇得到，不是爱是不是真的。人们放下的是遇见的成本，留下的是爱情。}
\begin{mechanism}\step 分开两种疑：疑遇见（机会、时间、能力）和疑爱情（不存在、不值得）\step 数据里的疑几乎全是前一种；婚、育、恋三样里恋爱被放下得最少\step 遇见被卡住时，嗑CP 让这份信保持温度，对多数人是补充不是替代\end{mechanism}
\begin{evidence}{数据}[已核实]八千多名单身受访者中 76.8\% 想谈恋爱（30 岁以下男性 85.9\%）；没脱单的前三原因：社交圈子单一 64.8\%、很难遇到三观契合的对象 48.4\%、不太擅长与异性交流 40.7\%。\sub{来源}{复旦发展研究院传播与国家治理研究中心、哔哩哔哩公共政策研究院，《中国青年网民社会心态调查报告（2024）》，复旦发展研究院网站 2025-01-20。}\sub{方法}{网络问卷 11253 人，其中单身者 8000 多人；另分析社交平台话题与帖子。}\sub{局限}{网民样本，偏年轻、偏城市；问的是想不想，不直接测信不信，所以要配合原因一起用。}\end{evidence}
\begin{evidence}{数据}[已核实]中国科学院心理研究所调查 55781 名在校大学生与 7366 名成年人：不愿脱单的人数比例低于不愿结婚和不愿生育子女的人数比例。\sub{来源}{陈祉妍、蒋建景、郭菲，《2024年成年人与在校大学生婚育观调查报告》，《中国国民心理健康发展报告（2023～2024）》，社会科学文献出版社 2025。}\sub{方法}{2024 年 3–6 月全国问卷。}\sub{局限}{测的是意愿，不是信念；新闻转述的“45.4\% 大学生认为拥有爱情不重要”只见于二手报道，未进稿件。}\end{evidence}
\begin{evidence}{学理}[已核实]类社会关系（Horton 和 Wohl，1956）\sub{核心主张}{观众与媒介人物之间单向的“隔着距离的亲密”；对大多数观众，它是日常社交生活的补充，日常关系里的预设在其中被展示和重申。}\sub{支撑}{机制第三步：嗑CP 让对爱情的预设被重申，而不是被取消。}\sub{边界}{同一篇写到：变成自主社交的替代时，才算病态；反方会引“随时可以抽身、几乎没有义务”。}\sub{口头版}{隔着屏幕的关系，对多数人是补充，不是替代。}\sub{出处}{Psychiatry 19(3): 215–229，doi:10.1080/00332747.1956.11023049。}\end{evidence}
\begin{evidence}{案例}[已核实]中国青年报采访的单身青年张钰（化名，28 岁）：婚恋负面消息常上热搜，让她“对感情提不起信心”，但她说“现在我一个人有点累了，想遇到对的人一起走，两人彼此信任，互相扶持”。\sub{代表性}{个案，用来说明“疑在上、信在下”的机制，不证明普遍。}\sub{反噬风险}{反方会引她同一篇里的“没必要冒着风险去试错”；正方用她的完整故事，最后一句是想遇到对的人。}\sub{出处}{杜园春，中国青年报 2024-05-30 第 3 版。}\end{evidence}
\closing{疑在上面，信在底下：拿掉“遇见”这层门槛，信还在；拿掉信，就没人嗑“双向奔赴”。}
\begin{clash}
\pair{想谈不等于信，想而不做正是疑。}{我方不只看想不想，看的是理由：没脱单的理由没有一条是爱情是假的。}
\pair{疑条件就是疑爱情：扛不住圈子窄的爱就是靠不住。}{靠不住说的是两个人之间的爱会不会变；圈子窄说的是两个人能不能遇见。门口挤，不等于门里没人。}
\pair{Horton 和 Wohl 说变成替代就是病态，一半人愿意只嗑不谈。}{那是两百三十多份的学生问卷；他们把多数人的情况叫补充，反方要证明多数是替代。}
\end{clash}
\end{fields}

\subsection{（被淘汰的正方候选论点，交给反驳库与对辩备用）}
\begin{enumerate}[leftmargin=1.8em]
\item 失望源于信：塌房、BE 让人难过，说明投入的是信。淘汰：可被反用成“一塌房就不信，这份信很脆”。
\item 代餐说明饿、爱情排在婚育前、社群确认、语言外溢：并入论点一、二，不单列。
\end{enumerate}

\consection{六、反方论点}
\consubsection{论点一：信在远处}
\begin{fields}
\claim{嗑CP 的爱永远在屏幕那头。人们相信爱情会发生在别人身上，却不相信它会、也不敢让它落到自己身上。}
\begin{mechanism}\step 嗑CP 的形式是旁观：那段爱，我进不去，也不用进\step 热潮越大，越多人把爱放在一个自己不必进场的位置\step 只对别人成立、对自己不成立的信，底下是对爱靠不靠得住的疑\end{mechanism}
\begin{evidence}{学理}[已核实]类社会关系（Horton 和 Wohl，1956）\sub{核心主张}{观众在这种关系里几乎不承担义务、努力与责任，随时可以抽身；互动单向，不能共同发展。}\sub{支撑}{机制第一步：嗑CP 的爱是一份不用付出、随时能走的爱。}\sub{边界}{同一篇说对大多数观众是补充；反方要用“人们自己的恋爱在退场”去论证补充正在变成替代。}\sub{口头版}{隔着屏幕的关系，不用付出，随时能走。}\sub{出处}{Psychiatry 19(3): 215–229。}\end{evidence}
\begin{evidence}{数据}[已核实]2905 名 18–26 岁未婚城市青年中，12.8\% 选择将来“不谈恋爱”，26.3\% 表示不确定。\sub{来源}{共青团中央中国特色社会主义理论体系研究中心课题组，《青年婚恋意愿调查：面对婚姻，年轻人在忧虑什么？》，光明日报 2021-10-08，央视网转载。}\sub{方法}{问卷调查加深度访谈。}\sub{局限}{对象是全体未婚城市青年，不是嗑CP 的人；只作背景，不能说“近四成不谈恋爱”。}\end{evidence}
\begin{evidence}{案例}[已核实]许冠文在采访中说，CP 粉里流传着一句话：“虽然我不谈恋爱，但是我要看我的CP谈恋爱。”他自己的疑问是：为什么她们会偏向于消费爱情的产品，而不是自己进行爱情的实践？\sub{代表性}{研究者对 CP 粉圈的概括，不是某一位受访者。}\sub{贴合度}{正好是“爱在别人那里，自己是‘虽然我不’”的结构。}\sub{出处}{红星新闻 2023-10-14 采访许冠文，凤凰网转载。}\end{evidence}
\begin{evidence}{补充数据}[已核实]人大本科生问卷：72.32\% 认为嗑CP 的幸福快乐程度不亚于谈恋爱，超过一半愿意只嗑CP 不谈恋爱。\sub{来源}{傅煜翔、迪丽娜尔，澎湃新闻·湃客 2023-12-31（中国人民大学人口与发展研究中心靳永爱责编）。}\sub{局限}{234 份有效答卷，便利样本；只在防守与对辩里用，说出样本量。}\end{evidence}
\closing{底下的是疑：爱情可以存在，但不在我这里，也不靠得住到值得我进场。}
\begin{clash}
\pair{旁观不等于怀疑，看球的人不上场也信足球。}{看球的人没说过“球场太危险，我只看”；嗑的人说了：现实里体验这种爱，很难不发疯。}
\pair{一成多不谈，那近九成呢？}{数据只是背景；主证据是嗑CP 的人自己的话。而且超过四分之一连自己会不会谈都不确定。}
\pair{那句话说“但是我要看”，想看就是信。}{信它会发生，但只在别人身上；对自己，是“虽然我不”。}
\end{clash}
\end{fields}

\consubsection{论点二：要的是保险}
\begin{fields}
\claim{人们嗑CP，是因为它是一份上了保险的爱：结局自己写好，不会解散，不会受伤。只有结局提前写好才敢投入，说明不信真实的爱靠得住。}
\begin{mechanism}\step 真实的爱必然带风险：会变、会散、会受伤\step 嗑CP 把风险提前拿掉：结局自己写，随时能退\step 当事人选它的理由正是“现实太险”，这就是对爱靠不靠得住的判断\end{mechanism}
\begin{evidence}{学理}[已核实]安全威胁（巴迪欧，《爱的多重奏》，2009）\sub{核心主张}{婚恋网站承诺“没有痛苦的完美爱情”（Get perfect love without suffering!），把风险提前算尽；巴迪欧称之为爱面对的“安全威胁”。}\sub{支撑}{机制第二、三步：先保结局再投入，是同一种逻辑。}\sub{边界}{巴迪欧批的是婚恋网站，不是嗑CP；只借“先保结局”的逻辑，不说他评论过 CP。}\sub{口头版}{先把风险算尽再去爱，巴迪欧叫它安全威胁。}\sub{出处}{Alain Badiou，Éloge de l'amour，2009；邓刚译，华东师范大学出版社 2012。}\end{evidence}
\begin{evidence}{补充学理}[有把握]液态之爱（鲍曼，2003）\sub{核心主张}{现代人渴望亲密与依靠，又害怕被关系长久拴住。}\sub{出处}{Zygmunt Bauman，Liquid Love，Polity，2003（出版社简介，未读正文）。}\end{evidence}
\begin{evidence}{数据}[已核实]不想恋爱的青年里，74.8\% 选择“一个人很好，谈恋爱很麻烦”。\sub{来源}{共青团中央课题组，光明日报 2021-10-08。}\sub{局限}{“麻烦”也可以被读成成本；同一调查的 30.5\% 选的是“不相信婚姻”，不是爱情，而且 73.4\% 是选此项者里女性的占比，不是女性里的比例，禁止误读。}\end{evidence}
\begin{evidence}{案例}[已核实]许冠文研究中的一位受访者：“如果在现实生活中去体验这种爱情，自己很难不发疯。”许冠文解释：“这种爱情带来的风险是非常大的。”记者问 CP 的倒推逻辑是否更安全，他答：“对。你不用担心他们会解散，因为你已经给他们画好了终点。”\sub{代表性}{研究者对 53 名深访对象的概括加一位受访者原话。}\sub{反噬风险}{同一采访里许冠文也说“我们对这种良好的亲密关系赋予了很高的价值”，正方会引前半句；反方要背出整句，强调后半句“每一步都充满风险”决定了人们怎么做。}\sub{出处}{红星新闻 2023-10-14，凤凰网转载。}\end{evidence}
\closing{选择的理由就是心态：因为不放心真实的爱，才去嗑一份结局写好的爱。}
\begin{clash}
\pair{怕风险说明把爱当真。}{当真，和相信它靠得住，是两件事；当真又不敢，正是中立定义里的怀疑。}
\pair{想要好结局是所有故事的共性，古人听戏也要团圆。}{古人听完戏也自己去成亲；今天的人说的是现实太险，只投入写好结局的那份。}
\pair{画好终点，是相信爱能走到终点。}{相信它在故事里走得到，不相信它在现实里自己走得到，所以才要亲手画。}
\end{clash}
\end{fields}

\subsection{（被淘汰的反方候选论点，交给反驳库与对辩备用）}
\begin{enumerate}[leftmargin=1.8em]
\item 爱情重要性下降：大学生 45.4\% 认为拥有爱情不重要。淘汰：只见于新闻转述，核实不了；“不重要”也不等于“怀疑”。
\item 越假越好嗑、现实退场、代餐、宁要角色不要真人、“又”相信爱情：并入论点一、二或进反驳库。
\end{enumerate}

\prosection{七、正方反驳库（针对反方全部可能论点）}
\begin{rebuttals}
\rebuttal{1}{信在远处：嗑CP 的爱在屏幕那头，人们不信它会落到自己身上。}{旁观不等于怀疑；还没遇到，不等于不信会遇到。}{打机制：反方要证明人们站在场外是因为觉得场上的东西是假的。复旦 2024 没脱单的理由前三是圈子、遇见、交流，没有一项是“不信爱情”。打论据：共青团那份问的是全体未婚青年，不是嗑CP 的人，而且说将来不谈的是一成多。}{“那为什么不去遇而去嗑？” → 两件事不冲突，嗑CP 不占用遇见的机会；圈子窄、时间少，是门口的事。}{最终}
\rebuttal{2}{要的是保险：画好终点、零风险，说明不信真实的爱靠得住。}{怕受伤，说明把爱当真；画好终点，是相信爱能走到终点。}{打定义：怕风险不等于不信爱，怕手术的人不是不信医学。打论据：许冠文同一句话的前半是“我们对这种良好的亲密关系赋予了很高的价值”；巴迪欧批的是婚恋网站，不是嗑CP。}{“当真又不去，不就是不信它靠得住？” → 不信的是自己能不能遇到对的人，不是爱本身；张钰最后说的是想遇到对的人。}{最终}
\rebuttal{3}{婚恋退潮：结婚登记从 2013 年的 1346.93 万对跌到 2024 年的 610.6 万对。}{结婚登记跟着条件走：二〇二五年全国通办后一年多了 65.7 万对，爱情观不会一年变一次。}{打自变量：结婚不等于爱情；中科院心理所 2024：不愿脱单的比例低于不愿结婚和不愿生育，人们放下的是婚育制度，留下的是爱情。}{“结婚又降了呢？” → 正好说明登记数随条件起落，拿它量不出信。}{常见}
\rebuttal{4}{爱情代餐：七成多认为嗑CP 快乐不亚于恋爱，一半愿只嗑不谈。}{那是两百三十多份的学生问卷；而且代餐说明饿。}{打论据：便利样本；打机制：人只给想要的东西找代餐，没人给自己不信的东西找代餐。许冠文本人说，只停留在“安抚毯”式替代性满足是不够的。}{“饿了只吃代餐？” → 饭馆离得远（遇见难），代餐让胃口还在。}{常见}
\rebuttal{5}{“我又相信爱情了”的“又”，说明本来不信。}{“又”说明信还能被唤回来；真正的不信，是怎么看都不动。}{打定义：怀疑的人看到甜会说“剧本”“炒作”；热潮里的人说的是“又信了”。Horton 和 Wohl：日常关系的预设会在这种关系里被重申。}{“靠别人才能信，是不是很脆？” → 信仰本来就需要被重申，宗教也要仪式。}{淘汰}
\rebuttal{6}{大学生 45.4\% 认为拥有爱情不重要。}{同一份报告的摘要说：不愿脱单的比例低于不愿结婚和不愿生育。}{打论据：45.4\% 只见于新闻转述；“不重要”是排序，不是“不信”。}{“近一半说不重要还不是疑？” → 请先出示原报告；排序低也不等于认为爱是假的。}{淘汰}
\rebuttal{7}{CP 粉宁要角色不要真人，逼真实情侣让路（2026 年 9 月《早春晴朗》事件）。}{极端个案，被广泛批评；而且他们执着的是剧里的专一，这份执念本身是信。}{打代表性：主流嗑CP 不攻击任何人；事件仍在发酵，细节不要在场上展开。}{“为什么宁要假的？” → 他们要的是剧里那份专一，不是不要真的。}{淘汰}
\rebuttal{8}{越假越好嗑：纸片人、虚构、BE 美学，被信的是永远不会被验证的爱。}{嗑虚构的人，嗑的仍是真心；虚构让标准不被现实打折，恰恰是守标准。}{打机制：虚构不是为了逃避，而是为了把理想完整地看一遍；标准会被带回现实。}{“守一个碰不到的标准是不是逃避？” → 标准带回了现实：宁缺毋滥。}{常见}
\end{rebuttals}

\consection{八、反方反驳库（针对正方全部可能论点）}
\begin{rebuttals}
\rebuttal{1}{嗑的是理想：CP 的标准是唯一、坚守、双向，说明信。}{标准高，是因为不用兑现；叶公对龙的标准也很高。}{打机制：想象里的标准不花成本；向往和相信之间隔着一个“敢”字。打自变量：浪漫信念量表测的是对自己恋爱的信念，不是对 CP 的期待。}{“标准被带回现实：宁缺毋滥” → 宁缺毋滥落到现实就是缺；拿标准拒绝一切，是疑不是信。}{最终}
\rebuttal{2}{疑的是条件：将近八成想谈恋爱，卡在圈子和遇见。}{想要不等于相信；想而不敢，正是怀疑的样子。}{打定义：中立定义里的怀疑包括“不信它靠得住”，一份扛不住圈子窄、扛不住负面热搜的爱就是靠不住。打案例：张钰原话“对感情提不起信心”“没必要冒着风险去试错”。}{“她最后说想遇到对的人” → 想遇到，但不去试：想在上面，不敢在底下。}{最终}
\rebuttal{3}{代餐说明饿。}{饿却只吃代餐，说明不信正餐吃得到。}{打机制：代餐的意义是替代；人大问卷超过一半愿意只嗑不谈（便利样本，说出样本量）。}{“饭馆远而已” → 远就不去了，说明不信那顿饭值得走这段路。}{常见}
\rebuttal{4}{失望源于信：塌房、BE 让人难过，说明投入了信。}{一次塌房就喊“再也不信爱情了”，这份信有多脆？}{打机制：信要靠别人不塌房来维持，就是不相信它自己站得住。}{“难过说明投入” → 投入的是旁观的情绪，不是自己的关系。}{淘汰}
\rebuttal{5}{爱情排在婚育前：中科院心理所调查，不愿脱单的少于不愿结婚、生育。}{最不被放弃，不等于被相信。}{打自变量：意愿不是信念；同一份报告的结论是青年婚恋意愿总体偏低。}{“至少没放弃爱情” → 没放弃，只是放在了屏幕那头。}{淘汰}
\rebuttal{6}{社群确认：嗑CP 是集体仪式，反复确认爱情存在。}{需要一群人天天确认，恰恰说明一个人信不住。}{打机制：确认的是“他们的爱是真的”，没人确认“我的也会是真的”。}{“宗教也要仪式” → 宗教仪式要你自己入场，嗑CP 不要。}{淘汰}
\rebuttal{7}{“我又相信爱情了”是信的宣言。}{“又”字说明之前不信；靠一对 CP 才能信，信在别人身上。}{打折射：题目看现象折射了什么，不看口号。}{“能被唤回说明底下是信” → 唤回的只是对他们的信，关掉屏幕还剩多少？}{常见}
\rebuttal{8}{嗑CP 让人学会爱、提高标准、练习共情。}{看一千场球学不会踢球；练习的前提是上场。}{打论据：没有证据显示嗑CP 的人更会谈恋爱；许冠文说现实的亲密关系“每一步都充满风险”。}{“至少提高了标准” → 标准越高越不敢上场，这就是保险逻辑。}{常见}
\end{rebuttals}

\section{九、持方优劣评估}
\begin{assessment}
\assess{定义空间}{中立定义里的怀疑包括“不信它靠得住、值得投入”，给反方留了“不敢算不信”的空间；信仰的“存在、值得”给正方留了内容空间。“更深”的表里结构两边可以对称使用。}{均衡}
\assess{论证义务}{双方都要证明自己的心态在底下、是另一种的来源。正方多一道坎：必须把“疑”锁在遇见层面，否则“疑条件就是疑爱情”会吞掉论点二；反方多一道坎：必须证明旁观是因为“不敢”，而不只是娱乐习惯。}{均衡}
\assess{论据可得性}{正方有两份大样本调查（复旦 2024 的将近八成想谈、中科院心理所五万多名大学生的意愿排序），都已核实；反方的核心证据是许冠文 53 人的质性访谈和一份 234 份的学生问卷，人口统计只能作背景。许冠文一句话前后两半各归一方。}{略偏正}
\assess{评委直觉}{没有评委信息，只作一般判断：“嗑CP 是爱情代餐”“年轻人不敢谈恋爱”在舆论里更常见，而“更深”这个词天然邀请“表面甜、底下苦”的读法，普通听众先倾向反方；正方要先克服“你们在美化逃避”的直觉。}{偏反}
\assess{赛制顺序}{正方先立论、最后结辩，正一结辩有临场位回应反一；反方先申论（反二），反三申论可以回应正三，反一结辩只能封路。正方收尾优势明显。}{略偏正}
\end{assessment}
\begin{fields}
\claim[结论]{正方略劣。劣势主要来自评委直觉：题目的“更深”和“代餐”叙事都顺着反方；正方的优势在论据和最后一发。}
\begin{procard}{劣势方打法}收窄战场：把全场拉到“疑的对象是什么”，一辩稿就把“疑遇见 / 疑爱情”分开，后场每个人都用“门口 / 门里”这一对词。判准之争：反复提醒判准有两半，反方只讲了形式。价值层：正一结辩的价值升华讲“等的人心里那把尺子”，把“美化逃避”的直觉翻成“守住标准”。顺序利用：正一最后一发用临场位接反一结辩的原话。\end{procard}
\begin{concard}{优势方要防的}反方最容易滑向“年轻人不相信爱情了”这种绝对说法，或把共青团数据说成“近四成不谈恋爱”，一被正方抓到就丢掉评委的信任。反方还要防被拖进比人数：主证据永远是嗑CP 的人自己的话。许冠文同一句话前半归正方，反方要主动说完整句。\end{concard}
\end{fields}

\section{十、口径表}
\prosubsection{10.1 正方口径表}
\begin{canon}{正方}
\canongroup{定义}
\canonrow{嗑CP}{以旁观者的身份，给虚构角色或者真实人物配对，想象并期待他们之间的浪漫关系，并从中得到情感满足。}{给角色或真人配对、期待他们在一起}{把嗑CP 说成“追求爱情”}{中立定义}
\canonrow{信仰 / 怀疑}{对爱情的信仰，是相信真正的爱情存在、值得期待和投入；怀疑，是不相信它存在，或者不相信它靠得住、值得投入。}{信，是信它存在、值得}{“喜欢爱情故事就算信”（定义杀）}{中立定义；汉典}
\canonrow{更深}{热潮里两种心态都有，比哪一种在底下：它是另一种的来源。}{谁在底下；谁是谁的来源}{“反方的疑根本不存在”}{中立定义}
\canongroup{判准}
\canonrow{中立判准}{哪一种心态更能解释热潮的核心特征：人们嗑的是什么样的爱，为什么用旁观的方式去爱，并且是另一种心态的来源。}{判准有两半：嗑什么，为什么这样嗑}{“只看嗑的内容”}{中立判准}
\canonrow{双方义务}{我方证明嗑CP 的动力来自相信真爱存在，犹豫疑的是遇见；对方证明嗑CP 的动力来自不信真爱靠得住。}{疑在上面，信在底下}{}{中立判准}
\canongroup{论点}
\canonrow{嗑的是理想}{人们嗑CP，嗑的是最严格的爱情标准：唯一、坚守、双向、好结局；肯为一个还没在现实里见过的标准投入热情，并拿它拒绝将就，这就是信。}{嗑出来的是一把尺子；标准没有降，是升了}{“嗑CP 让人学会爱”；“双向奔赴是 CP 圈的词”；“怀疑的人不嗑CP”}{论点一}
\canonrow{疑的是条件}{热潮里的犹豫是真的，但它疑的是遇不遇得到，不是爱是不是真的；人们放下的是遇见的成本，留下的是爱情。}{卡在门口，不是门里没有人}{“年轻人都想谈恋爱”；“不谈恋爱都是因为没机会”；“结婚登记回升说明信爱情”}{论点二}
\canongroup{数据}
\canonrow{将近八成想谈}{复旦发展研究院 2024：八千多名单身受访者里，76.8\% 想谈恋爱。}{将近八成单身青年想谈恋爱}{“年轻人都想谈”；“超过八成”}{复旦发展研究院等 2025-01}
\canonrow{没脱单的原因}{社交圈子单一 64.8\%，很难遇到三观契合的对象 48.4\%，不太擅长与异性交流 40.7\%。}{排第一的是圈子窄，六成多的人选了}{“没有人选不信爱情”（问卷没有这一项，只说“三条里没有一条是爱情是假的”）}{同上}
\canonrow{婚育恋排序}{中科院心理所 2024 调查 55781 名大学生：不愿脱单的人数比例低于不愿结婚和不愿生育。}{五万五千多名大学生；婚、育、恋三样里，恋爱被放下得最少}{“大学生都想恋爱”；引 45.4\%}{陈祉妍等 2025}
\canonrow{恋爱期待}{中国青年报社会调查中心：58.5\% 希望互相欣赏、变得更加自信，58.1\% 期待彼此包容、更有安全感。}{将近六成希望互相欣赏}{报样本量（文中未公开）}{中国青年报 2024-05-30}
\canonrow{许冠文研究}{许冠文 2023 深访 53 名女性 CP 粉：CP 更像是一个将她们的期待具象化的载体。}{CP 是把期待具象化的载体}{“论文证明了嗑CP 的人信爱情”}{《妇女研究论丛》2023(1)；红星新闻 2023-10-14}
\canongroup{学理}
\canonrow{浪漫信念}{Sprecher 和 Metts 1989 年的浪漫信念量表：唯一、爱能找到出路、把对方理想化、一见钟情。}{CP 剧把这张量表演了一遍}{“量表证明 CP 粉更浪漫”}{JSPR 1989}
\canonrow{类社会关系}{Horton 和 Wohl 1956：对大多数观众，隔着屏幕的关系是日常生活的补充，日常关系里的预设在其中被重新确认。}{补充，不是替代}{“他们说嗑CP 很健康”}{Psychiatry 1956}
\canongroup{案例}
\canonrow{双向奔赴}{“双向奔赴”是《咬文嚼字》2023 年十大流行语，释义是“相互爱慕、相互亲近的美好愿望”。}{一个说爱情理想的词成了年度流行语}{说它来自 CP 圈；展开讲外交用法}{文汇报 2023-12-04}
\canonrow{张钰}{中国青年报采访的张钰（化名，28 岁）：负面热搜让她对感情提不起信心，但她想遇到对的人一起走，两人彼此信任。}{卡住她的是门口}{只引前半句或只引后半句}{中国青年报 2024-05-30}
\canongroup{分工与全场问题}
\canonrow{申论分工}{正二申论：拆反方立论主干，推进论点一（恋爱期待将近六成、宁缺毋滥）。正三申论：处理反方最强点，推进论点二（婚育恋排序、Horton 和 Wohl 补充）。}{二辩拆立论，三辩打最强点}{两篇讲同一件事}{赛制文件}
\canonrow{全场问题}{如果人们不信爱情，为什么偏偏嗑“唯一”和“双向奔赴”，而不嗑一对互相算计的人？}{不信爱情，为什么嗑专一}{}{一辩稿}
\end{canon}

\consubsection{10.2 反方口径表}
\begin{canon}{反方}
\canongroup{定义}
\canonrow{嗑CP}{以旁观者的身份，给虚构角色或者真实人物配对，想象并期待他们之间的浪漫关系，并从中得到情感满足。}{给角色或真人配对、隔着屏幕期待}{把嗑CP 说成“逃避现实”“病态”}{中立定义}
\canonrow{信仰 / 怀疑}{对爱情的信仰，是相信真正的爱情存在、值得期待和投入；怀疑，是不相信它存在，或者不相信它靠得住、值得投入。}{疑，是不信它靠得住、值得投入}{“不去谈恋爱就是不信”（定义杀）}{中立定义；汉典}
\canonrow{更深}{热潮里两种心态都有，比哪一种在底下：它是另一种的来源。}{谁在底下；谁是谁的来源}{“正方的信根本不存在”}{中立定义}
\canongroup{判准}
\canonrow{中立判准}{哪一种心态更能解释热潮的核心特征：人们嗑的是什么样的爱，为什么用旁观的方式去爱，并且是另一种心态的来源。}{判准有两半：嗑什么，为什么这样嗑}{“只看人们谈不谈恋爱”}{中立判准}
\canonrow{双方义务}{我方证明嗑CP 的动力来自不信真爱靠得住，热潮里的甜只在安全距离内成立；对方证明嗑CP 的动力来自相信真爱存在。}{信在上面，疑在底下}{}{中立判准}
\canongroup{论点}
\canonrow{信在远处}{嗑CP 的爱永远在屏幕那头；人们相信爱情会发生在别人身上，却不相信它会、也不敢让它落到自己身上。}{爱在别人那里；隔着屏幕的信}{“年轻人不相信爱情了”；“嗑CP 的人都不谈恋爱”；“嗑CP 是病态”}{论点一}
\canonrow{要的是保险}{人们嗑CP，是因为它是一份上了保险的爱：结局自己写好，不会解散，不会受伤；只有结局提前写好才敢投入，说明不信真实的爱靠得住。}{先保结局再投入；零风险的爱}{“怕风险的人都不信爱”}{论点二}
\canongroup{数据}
\canonrow{将来不谈 / 不确定}{共青团中央课题组 2021 调查 2905 名 18–26 岁未婚城市青年：12.8\% 选择将来不谈恋爱，26.3\% 不确定。}{两千九百多名未婚城市青年里，一成多说不谈，超过四分之一说不确定}{“近四成青年不谈恋爱”}{光明日报 2021-10-08}
\canonrow{嫌麻烦}{同一调查：不想恋爱的青年里，74.8\% 选“一个人很好，谈恋爱很麻烦”。}{七成半选了“一个人很好，谈恋爱很麻烦”}{“七成半的青年不想恋爱”；引 30.5\% 说成不信爱情；“女性 73.4\% 不相信婚姻”}{同上}
\canonrow{只嗑不谈}{人大本科生问卷 234 份：72.32\% 认为嗑CP 的快乐不亚于谈恋爱，超过一半愿意只嗑CP 不谈恋爱。}{一份两百三十多份的学生问卷里，七成多觉得嗑CP 不比恋爱差}{不报样本量；“一半年轻人只嗑不谈”}{澎湃新闻·湃客 2023-12-31}
\canonrow{许冠文研究}{许冠文 2023 深访 53 名女性 CP 粉；他说现实中的亲密关系“每一步都充满风险”。}{研究者自己说每一步都有风险}{只引后半句被正方点破时不认前半句}{《妇女研究论丛》2023(1)；红星新闻 2023-10-14}
\canongroup{学理}
\canonrow{类社会关系}{Horton 和 Wohl 1956：观众几乎不承担义务和努力，随时可以抽身；这种关系是单向的，不会一起往前走。}{不用付出，随时能走}{“他们说这种关系是病态”}{Psychiatry 1956}
\canonrow{安全威胁}{巴迪欧《爱的多重奏》：婚恋网站承诺“没有痛苦的完美爱情”，他叫它爱面对的“安全威胁”。}{先把风险算尽，再去爱}{“巴迪欧批评过嗑CP”}{Badiou 2009}
\canongroup{案例}
\canonrow{虽然我不}{许冠文说，CP 粉里流传着一句话：“虽然我不谈恋爱，但是我要看我的CP谈恋爱。”}{爱在别人那里，落到自己是“虽然我不”}{“一位受访者说”（归属错）}{红星新闻 2023-10-14}
\canonrow{画好终点}{一位受访者：在现实里体验这种爱情，自己很难不发疯；许冠文：你不用担心他们会解散，因为你已经给他们画好了终点。}{结局写好才敢投入}{把“画好终点”说成受访者原话}{同上}
\canongroup{分工与全场问题}
\canonrow{申论分工}{反二申论：拆正方立论主干，推进论点一（Horton 和 Wohl：不用付出、随时抽身）。反三申论：回应正三，处理正方最强点，推进论点二（巴迪欧安全威胁、七成半嫌麻烦）。}{二辩拆立论，三辩打最强点}{两篇讲同一件事}{赛制文件}
\canonrow{全场问题}{如果人们真信爱情，为什么要把它放在一段自己碰不到、结局早写好的关系里？}{真信，为什么只敢放在屏幕那头}{}{一辩稿}
\end{canon}

\section{十一、论据出处清单}
\begin{sourcelist}
\source{1}{单身青年想谈恋爱：8000 多名单身受访者中 76.8\%；30 岁以下男性 85.9\%}{复旦发展研究院传播与国家治理研究中心、哔哩哔哩公共政策研究院，《中国青年网民社会心态调查报告（2024）》，2025，\url{https://fddi.fudan.edu.cn/ef/f8/c18986a716792/page.htm}}{超过四分之三（76.8\%）的年轻人表示自己“想谈恋爱”}{2026-09-30}{已核实}{正一立论；速查}{}
\source{2}{没脱单原因：社交圈子单一 64.8\%，难遇三观契合 48.4\%，不擅长与异性交流 40.7\%}{复旦发展研究院等，《中国青年网民社会心态调查报告（2024）》，2025，\url{https://fddi.fudan.edu.cn/ef/f8/c18986a716792/page.htm}}{“社交圈子单一，很少有机会认识异性”（64.8\%）}{2026-09-30}{已核实}{正一立论；正二质询；速查}{}
\source{3}{婚育恋意愿排序：55781 名大学生、7366 名成年人，不愿脱单比例低于不愿结婚、不愿生育}{陈祉妍、蒋建景、郭菲，《2024年成年人与在校大学生婚育观调查报告》，《中国国民心理健康发展报告（2023～2024）》，社会科学文献出版社 2025，\url{https://cpms.ssap.com.cn/skwx_cmr/initDatabaseDetail?contentId=8225485&siteId=23&contentType=literature}}{不愿脱单（找对象）的人数比例低于不愿结婚和不愿生育子女的人数比例}{2026-09-30}{已核实}{正三申论；速查}{}
\source{4}{大学生 45.4\% 认为拥有爱情不重要}{腾讯新闻 2025-06-20 转述上条报告}{}{}{有把握}{备用（未进稿件）}{查蓝皮书纸本或皮书数据库全文，确认原句与分母}
\source{5}{恋爱期待：58.5\% 希望互相欣赏、变得更加自信；58.1\% 期待彼此包容、更有安全感}{中国青年报社会调查中心，杜园春《年轻人婚恋更看重情感质量》，中国青年报 2024-05-30 第 3 版，\url{http://zqb.cyol.com/html/2024-05/30/nw.D110000zgqnb_20240530_2-03.htm}}{58.5\%的受访青年希望是互相欣赏，变得更加自信}{2026-09-30}{已核实}{正二申论；速查}{}
\source{6}{张钰（化名）个案}{杜园春，中国青年报 2024-05-30，\url{http://zqb.cyol.com/html/2024-05/30/nw.D110000zgqnb_20240530_2-03.htm}}{想遇到对的人一起走，两人彼此信任，互相扶持}{2026-09-30}{已核实}{正一立论；反方反驳库}{}
\source{7}{周亚玲（化名）个案}{杜园春，中国青年报 2024-05-30，\url{http://zqb.cyol.com/html/2024-05/30/nw.D110000zgqnb_20240530_2-03.htm}}{不急不躁，宁缺毋滥}{2026-09-30}{已核实}{正二申论}{}
\source{8}{2905 名未婚城市青年：12.8\% 将来不谈恋爱，26.3\% 不确定}{共青团中央中国特色社会主义理论体系研究中心课题组，《青年婚恋意愿调查：面对婚姻，年轻人在忧虑什么？》，光明日报 2021-10-08，央视网转载 \url{https://news.cctv.com/2021/10/08/ARTIzxZbsSpcvsDH8RUqQoXk211008.shtml}}{12.8\%的青年选择“不谈恋爱”}{2026-09-30}{已核实}{反一立论；速查}{}
\source{9}{不想恋爱原因：74.8\% 选“一个人很好，谈恋爱很麻烦”}{共青团中央课题组，光明日报 2021-10-08，\url{https://news.cctv.com/2021/10/08/ARTIzxZbsSpcvsDH8RUqQoXk211008.shtml}}{选择“一个人很好，谈恋爱很麻烦”的比例为74.8\%}{2026-09-30}{已核实}{反三申论；速查}{}
\source{10}{30.5\% 选“不相信婚姻”，其中女性占 73.4\%}{共青团中央课题组，光明日报 2021-10-08，\url{https://news.cctv.com/2021/10/08/ARTIzxZbsSpcvsDH8RUqQoXk211008.shtml}}{青年选择“不相信婚姻”的占30.5\%，其中女性占73.4\%}{2026-09-30}{已核实}{备用（禁止误读）}{}
\source{11}{许冠文、张慧深访 53 名女性 CP 粉；“CP 更像是一个将她们的期待具象化的载体”}{许冠文、张慧，《“嗑CP”：青年女粉丝的创造性情感体验》，《妇女研究论丛》2023 年第 1 期；红星新闻采访 2023-10-14，凤凰网转载 \url{https://news.ifeng.com/c/8TsckB0EYgh}}{CP更像是一个将她们的期待具象化的载体}{2026-09-30}{已核实}{正一立论；速查}{}
\source{12}{许冠文采访：CP 粉流行语“虽然我不谈恋爱，但是我要看我的CP谈恋爱”}{红星新闻 2023-10-14 采访许冠文，凤凰网转载 \url{https://news.ifeng.com/c/8TsckB0EYgh}}{（CP粉丝中）有一句话是“虽然我不谈恋爱，但是我要看我的CP谈恋爱”}{2026-09-30}{已核实}{反一立论；速查}{}
\source{13}{许冠文采访：受访者“很难不发疯”；“风险是非常大的”；“画好了终点”}{红星新闻 2023-10-14 采访许冠文，凤凰网转载 \url{https://news.ifeng.com/c/8TsckB0EYgh}}{你不用担心他们会解散，因为你已经给他们画好了终点}{2026-09-30}{已核实}{反一立论；速查}{}
\source{14}{许冠文采访：“对这种良好的亲密关系赋予了很高的价值”/“每一步都充满风险”}{红星新闻 2023-10-14 采访许冠文，凤凰网转载 \url{https://news.ifeng.com/c/8TsckB0EYgh}}{现实中的亲密关系实践，无时无刻不在提醒我们，每一步都充满风险}{2026-09-30}{已核实}{双方对辩与质询}{}
\source{15}{人大本科生问卷 234 份：72.32\% 认为嗑CP 快乐不亚于恋爱；超过一半愿只嗑不谈；23.71\% 不想谈恋爱}{傅煜翔、迪丽娜尔，《为什么许多大学生喜欢磕CP，自己却不恋爱？》，澎湃新闻·湃客 2023-12-31，\url{https://www.thepaper.cn/newsDetail_forward_25836154}}{72.32\%的人认为磕CP的幸福快乐程度不亚于谈恋爱}{2026-09-30}{已核实}{反方防守与对辩；速查}{}
\source{16}{FUNJI《嗑CP心态调查报告》4521 份，99.2\% 女性，单身超九成}{经济观察网 李佩珊 2022-08-09 转述}{}{}{有把握}{备用（未进稿件）}{找 FUNJI 原发布（微博或公众号）核对}
\source{17}{结婚登记：2013 年 1346.93 万对；2024 年 610.6 万对（比上年下降 20.5\%）；2025 年 676.3 万对（增加 65.7 万对）}{民政部统计公报与 2026-02-11 发布；经观察者网 2026-02-11、经济观察网 2024-08-23 转述}{}{}{有把握}{正方反驳库；正方对辩备用}{民政部官网（本环境无法访问）《2024年民政事业发展统计公报》与 2025 年四季度婚姻登记数据}
\source{18}{“双向奔赴”入选《咬文嚼字》2023 年十大流行语}{《咬文嚼字》编辑部，文汇报客户端 2023-12-04，中国作家网转载 \url{https://www.chinawriter.com.cn/n1/2023/1204/c403994-40131507.html}}{表达了人们相互爱慕、相互亲近的美好愿望}{2026-09-30}{已核实}{正一立论；速查}{}
\source{19}{浪漫信念量表四条信念，730 名本科生}{Sprecher, S. \& Metts, S. (1989). Development of the Romantic Beliefs Scale. Journal of Social and Personal Relationships 6(4): 387–411. doi:10.1177/0265407589064001}{the four beliefs comprising the scale: Love Finds a Way, One and Only, Idealization, and Love at First Sight}{2026-09-30}{已核实}{正一立论；速查}{}
\source{20}{类社会关系：补充与重申；无义务、随时抽身}{Horton, D. \& Wohl, R. R. (1956). Mass Communication and Para-Social Interaction. Psychiatry 19(3): 215–229. doi:10.1080/00332747.1956.11023049（全文见 participations.org 转载）}{He is free to withdraw at any moment.}{2026-09-30}{已核实}{正三申论；反二申论；速查}{}
\source{21}{巴迪欧：安全威胁，“没有痛苦的完美爱情”}{Badiou, A. (2009). Éloge de l'amour；英译 In Praise of Love, 2012；中译《爱的多重奏》，邓刚译，华东师范大学出版社 2012，\url{https://book.douban.com/subject/19964304/}}{Get perfect love without suffering!}{2026-09-30}{已核实}{反三申论；速查}{}
\source{22}{鲍曼《液态之爱》：渴望亲密又怕被拴住}{Bauman, Z. (2003). Liquid Love. Polity（出版社简介）}{}{}{有把握}{备赛文档}{读正文第一章，找原句}
\source{23}{信仰的词典义}{汉典“信仰”条（2026 年访问），\url{https://zdic.net/hans/信仰}；《现代汉语词典》“拿来作为自己行动的榜样或指南”}{对某种主张、主义、宗教或某人极其相信和尊敬}{2026-09-30}{已核实}{定义}{}
\source{24}{嗑 / 磕 的词典义；现汉未收 CP}{水寒说语文，《组CP，嗑CP，炒CP及其他》，网易 2022-10-03}{}{}{有把握}{定义}{查《现代汉语词典》第 7 版“嗑”“磕”条}
\source{25}{《早春晴朗》CP 粉要求井柏然与刘雯分手“为剧中角色让路”}{东方日报（马来西亚）2026-09-27，\url{https://www.orientaldaily.com.my/news/entertainment/2026/09/27/850789}}{涌入刘雯的社群留言要求她分手、“为剧中角色让路”}{2026-09-30}{已核实}{正方反驳库（事件仍在发酵）}{}
\end{sourcelist}

\section{十二、备赛建议}
\begin{fields}
\begin{timeline}\when{第 1–2 天}{全队对齐中立定义与判准的“两半”，背出本方全场问题；双方都要能背出许冠文那句整话（很高的价值 / 每一步都充满风险）。}\when{第 3–4 天}{一辩背立论；二辩、三辩按分工练申论套件，每篇练遍全部模块组合；被质询预案对练。}\when{第 5–6 天}{三场对辩各练两轮，质询按 90 秒双边计时练，记录是否打断。}\when{第 7 天}{完整模辩一场，按场上记录卡组装申论和结辩；复盘改口径表。}\end{timeline}
\field{模辩重点检验}{\sub{一}{正方能不能把每次“为什么不去谈”都拉回“门口 / 门里”。}\sub{二}{反方能不能在“怕说明当真”之后说出“当真又不敢，就是不信它靠得住”。}\sub{三}{质询有没有频繁打断，被质询有没有拖时间。}}
\begin{checklist}\item 我方定义：嗑CP、信仰 / 怀疑、更深\item 我方判准：两半（嗑什么、为什么这样嗑）+ 谁是来源\item 论点标签与一句话\item 许冠文那句话的完整版和归属\item 本方全场问题\item 禁说清单（口径表“禁止”一列）\end{checklist}
\end{fields}
\end{document}
